%% =====================================================================
\section{Engine and Benchmark}
\label{sec:engine}

All methods run on one engine~\cite{rubikrepo}. Everything below is generic over cuboid dimensions
$(d_0,d_1,d_2)$, so the cube $(N,N,N)$ and the Square $(N,N,1)$ come from the same code.

\paragraph{Moves from geometry.}
Each sticker stores a doubled, centred cubie coordinate and an outward normal. A move rotates every sticker of one
slab by an integer rotation matrix; looking up the rotated (position, normal) pair gives a permutation $\pi$, and
applying the move is the gather $s\leftarrow s[\pi]$, $O(6N^2)$ time. A slab whose cross-section is not square, such
as one row of the Square, only gets 180$^\circ$ turns, and an axis of length 1 gets none. Unit tests check that every
move has order 4 (or 2) and that each move followed by its inverse is the identity, for every tested $N$.

\paragraph{Canonical form.}
Moves on the same axis commute. We canonicalise a sequence by merging each same-axis run per layer and sorting by
layer, and we use the same order to prune search: after a move on axis $a$, layer $\ell$, no move on axis $a$ with
layer $\le\ell$ is allowed. This removes immediate repeats and every redundant ordering of commuting moves, and it is
the word model of the counting bound (\S\ref{sec:bound}).

\paragraph{Clusters with $N$-invariant signatures.}
Union-find over all move permutations gives the sticker orbits; orbits sharing cubies are merged into cubie clusters.
Inside a cluster, stickers are sorted by face and by the signs and relative magnitudes of their in-face coordinates.
This order does not depend on $N$, so the moves acting on a cluster induce the same local permutations at every $N$.
We hash the local generator set and the solved pattern into a \emph{signature}; tests confirm that, for example, every
$\times$-center cluster of the \Cu{5}, \Cu{7} and \Cu{9} has the same signature. A model trained for one cluster type
therefore applies at any size.

\paragraph{Projections.}
Restricting attention to a subset of stickers gives a \emph{projected puzzle}. When it is small we enumerate it by
breadth-first search from the solved pattern; the resulting table gives exact distances for that subset and doubles
as an exact membership test.

\paragraph{Solvability checker.}
A state is accepted only if it passes: global and per-cluster colour counts; membership in the exact projection table
for every cluster small enough to enumerate (all Square clusters); for cubes, valid corner, middle-edge and wing pieces,
corner twist $\equiv 0 \pmod 3$ and, for odd $N$ after normalising orientation by the fixed centers, even middle-edge
flip and equal corner and middle-edge permutation parity; and for the Square, joint consistency of each cross cluster
with the border clusters it shares moves with. These are necessary conditions. Property tests confirm that the checker
accepts scrambled states and rejects a single corner twist, a single middle-edge flip, a swap of two corners on odd
cubes and wrong colour counts.

\paragraph{Harness.}
Every solver implements one interface and returns engine moves. The harness re-applies the moves, decides
\emph{solved} itself, and writes one row per instance with lengths in all three metrics, wall time, nodes, memory,
seed, hardware and git commit. A crash or time-out counts as a failure and stays in the denominator. External programs
(KMCoders' solver, dwalton76, \texttt{nxopt}) run as subprocesses behind the same interface: their answers are parsed
into engine moves (wide turns expanded into single-slice turns) and re-verified like any other method's.

\paragraph{Shared beam search.}
\Mfour, \FiveA and \FiveB share one beam search. Per step it expands the beam in chunks, scores children with the
heuristic, and keeps the $W$ best after removing duplicates by a random 64-bit linear hash (also across the last three
steps). Only (score, hash, parent, action) is stored per child; the next beam is rebuilt from parents. Beam states are
packed at 3 bits per sticker and can live on the device, in pinned host memory or on disk. The search accepts
\emph{macro actions}: fixed move sequences applied as one step and expanded into ordinary moves in the output. A
heuristic that can score all children of a parent in one pass is used that way, as in Q* search~\cite{agostinelli2024qstar}.
One step of a width-$W$ search with $A$ actions costs $O(WA)$ heuristic evaluations and $O(WA\cdot N^2)$ sticker work.

\paragraph{Pure 3-cycles for any cluster.}
\Mtwo, \Mthree and \FiveC need, for every center and wing cluster of every $N$, move sequences that 3-cycle three of its
pieces and touch nothing else. We find them automatically. For centers we enumerate commutators $[A, B\,C\,B']$ with
$A$ and $C$ inner-slice quarter turns on one axis and $B$ an outer face turn; for wings, $[X, Y\,Z\,Y']$ with $X$ an
inner slice and $Y, Z$ outer faces. A candidate is kept if, applied to the sticker permutation, it moves exactly three
cubies of one cluster. A conjugate $s\,c\,s'$ of a pure 3-cycle $c$ is again pure, so breadth-first search over setup
moves $s$ (up to four) extends a few base commutators to every one of the $24\cdot23\cdot22/3=4{,}048$ 3-cycles of a
cluster, each with a short move sequence.
